% !TEX root = chapter-standalone.tex


\devel{
\section{Further theory}
\label{sec:functors-further-theory}

\subsection{Full and faithful functors}

\begin{ctdefinition}[Full and faithful functors]
    \label{def:functorfullfaith}
    \SYNDEF{full functor}
    \SYNDEF{faithful functor}
    \label{def:full-functor}
    \label{def:faithful-functor}
    A \SY{functor}~$\funa\colon \CatC \fto \CatD$ is \emph{full} (respectively \emph{faithful}) if for each pair of objects~$\Obja,\Objb\setin \CatC$, the function
    \begin{equation}
        \funa \colon \HomSet{\CatC}{\Obja}{\Objb} \to \HomSet{\CatD}{\funa(\Obja)}{\funa(\Objb)}
    \end{equation}
    is \SY{surjective} (respectively \SY{injective}).
\end{ctdefinition}

\begin{marginfigure}
    \centering
    \includesag{095_full_faithful_1}
    \caption{}
    \label{fig:ex_full_faithful_1}
    \todomistake{\alphubel: This figure is wrong}
\end{marginfigure}
\begin{example}
    Let~\CatC be the category depicted in \cref{fig:ex_full_faithful_1}.
    Let~$\funa\colon \CatC \fto \CatC$ be the \SY{endofunctor} which maps object~$\Obja\setin \Obof{\CatC}$ to~$\Obja$ and object~$\Objb\setin \Obof{\CatC}$ to~$\Obja$.
    This \SY{functor} is \SY{full} and \SY{faithful}.
    Note that the map of objects and the map of morphisms are neither \SY{surjective} nor \SY{injective}.
\end{example}

\begin{example}
    Consider a \SY{functor} which maps a category \CatC to its \SY{pre-order} structure \CatD, that is a category with the same objects as \CatC and a \emph{unique} morphism from~$\Obja\setin \Obof\CatD$ to~$\Objb\setin \Obof\CatD$ if and only if there is at least one morphism from~$\Obja$ to~$\Objb$ in \CatC.
    This \SY{functor} is \SY{full} by construction, and it is \SY{faithful} if and only if \CatC is a \SY{pre-order}.
    \todojira{136}{maybe add pic.
        OK, so what?
        Examples?
        Why is it useful?
    }
\end{example}

\subsection{Embeddings}
\publictodomessage

\linkvideo{spring2021-functors:semi-and-fun:embedding-fun} % Embedding functors
\begin{ctdefinition}[Embedding functor]
    \label{def:embedding-functor}
    \SYNDEF{embedding functor}

    A \SY{functor}~$\funa\colon \CatC \fto \CatD$ is an \emph{embedding} if it is faithful and~$\funaob$ is \SY{injective}.
\end{ctdefinition}

\todotextjira{133}{Write section on embeddings and merge duplicates.}


\subsection{Forgetful functors}
\label{sec:forgetful-functors}
\linkvideo{spring2021-functors:semi-and-fun:forgetful-fun} % Forgetful functors

\todotextjira{411}{\alphubel: @JL: Give more examples of forgetful \SY{functors}, and more explanation/intuition.
}

\begin{example}
    The \SY{functor}~$\Pos \fto \Set$ is a forgetful \SY{functor}.
    This \SY{functor} maps every \SY{poset} to the set which has the same elements, but no notion of order.
    Furthermore, it maps each \SY{monotone map} between \SY{posets} to the corresponding function between sets.
    This is a forgetful \SY{functor} in the sense that it forgets the notion of order and \SY{monotone maps}.
\end{example}


\section{Roles of functors}

In this functor part of the book, perhaps we can organize the applications part into several groups, bases on an operational understanding of the different ways that functors can be useful. 

Here are notes from the discussion on 29.06.2026:

The possible jobs of a functor:
\begin{enumerate}
\item embedding-like (also deifnitng structrure of A from a map A → B where B has structure). the functor defines the structure. you choose nay function, you assert that it must be fucntor

\item compiling-like, projection like. You need to find a functor (loses infomration, but preserves the “Structure”)

example: a compiler source code → machine code. 

example: the deployment computation grah → resource graph

\item there is a function, and a property of it is encoded by the enforcing of it being a functor

example:  monotone functions 

\item the functor describes a lifting/generaliation -  a recipe / construatction

\item the functors encodes constraints that something must respect. (the functor is there before the categroy). spivak: the functor encodes integrity constrinats on the data - you REJECT data that is not compatible; this is generalize in a sense to enocoding diagrams via functors 

\item the search for a functor: the solution of a problem consists of specifying a functor
\end{enumerate}

example: berg


}

